\section{Einleitung}

\subsection{Motivation und Problemstellung}

\ac{E2E}-Tests prüfen eine Anwendung aus Nutzersicht und sind daher gerade für komplexe und interaktive Oberflächen wertvoll. Ihre Erstellung ist jedoch aufwendig, da jeder Testfall die Bedienschritte und die zu prüfenden Zustände einzeln nachbilden muss. In der Praxis führt das dazu, dass \ac{E2E}-Tests oft erst spät oder nur für einen Teil der vorgesehenen Abläufe umgesetzt werden. \acp{LLM} haben in den letzten Jahren in der Softwareentwicklung stark an Bedeutung gewonnen und werden zunehmend eingesetzt, um diesen Aufwand zu verringern und Tests direkt aus einer fachlichen Beschreibung zu generieren (\cite[S.~1]{hou_large_2024}; \cite[S.~1]{ferreira_acceptance_2025}). Für die Testerstellung ist das ein attraktiver Weg, da Use Cases oder User Storys in vielen Projekten ohnehin vorliegen und unmittelbar als Eingabe dienen können.

Was sich schnell erzeugen lässt, muss allerdings auch überprüft werden. Durch \acp{LLM} sinkt der Aufwand für die Testerstellung, der Aufwand für die Bewertung hingegen nicht. Ob ein generierter Test tatsächlich prüft, was die Anforderung verlangt, lässt sich durch sein Ausführungsergebnis nicht erkennen. Je nach Design kann er fehlerhaften Code fälschlich bestätigen, statt ihn aufzudecken, und so ein trügerisches Sicherheitsgefühl erzeugen \parencite[S.~5]{mathews_design_2024}. Die Herausforderung ist also nicht mehr, eine große Menge an Tests zu erzeugen, sondern ihre Qualität sicherzustellen.

Ohne Informationen über die tatsächlich vorhandenen Elemente einer Anwendung kann ein \ac{LLM} nur plausibel erscheinende Bezeichnungen raten. Ein generierter Test kann dann syntaktisch korrekt sein, aber trotzdem bereits an der ersten Interaktion scheitern, weil das angesprochene Element in der Anwendung gar nicht existiert.

Bei webbasierten \ac{GIS}-Anwendungen verschärfen sich die Probleme zusätzlich. Geoportale von Behörden sowie Umwelt- und Klimadatenportale stellen Geodaten im Browser bereit und werden inzwischen von einem breiten Publikum genutzt, sodass Fehler in der Kartendarstellung unmittelbar sichtbar werden. Gleichzeitig sind solche Anwendungen schwer automatisiert zu testen, da die Karte technisch bedingt nicht wie gewöhnliche Oberflächenelemente ansprechbar ist (vgl. Abschnitt~2.4). Gerade die Interaktion mit der Karte bildet aber den fachlichen Kern solcher WebGIS-Anwendungen und müsste daher von einem aussagekräftigen Test erfasst werden. Wie viel und welche Art von \ac{UI}-Kontext ein \ac{LLM} für eine zuverlässige Testgenerierung benötigt, ist allgemein wenig untersucht (vgl. Abschnitt~2.5). Für WebGIS-Anwendungen mit Karteninteraktionen existiert dazu bislang keine Untersuchung.

Diese Arbeit entsteht in Zusammenarbeit mit der con terra GmbH, in deren Projekten webbasierte \ac{GIS}-Anwendungen entwickelt und getestet werden. Aus dieser Praxis heraus stellt sich die Frage, wie viele Informationen über die Oberfläche einem \ac{LLM} zur Testgenerierung bereitgestellt werden müssen und in welcher Form diese Informationen den größten Nutzen bringen.

\subsection{Ziele}

Diese Arbeit untersucht, welchen Einfluss der im Prompt bereitgestellte \ac{UI}-Kontext auf die Qualität \ac{LLM}-generierter Playwright-\ac{E2E}-Tests aus Use Cases für WebGIS-Anwendungen auf Basis von \ac{OPT} hat. \ac{OPT} ist ein Open-Source-Framework zur Entwicklung webbasierter GeoIT-Anwendungen. Als Qualität wird dabei sowohl die Ausführbarkeit der Tests als auch die inhaltliche Korrektheit der Prüfung verstanden.

Daraus ergeben sich drei Teilfragen:

\begin{enumerate}
    \item Führt ein umfangreicherer \ac{UI}-Kontext zu einer höheren Ausführbarkeit der generierten Tests?
    \item Wirken sich unterschiedliche Formen des Kontexts auf unterschiedliche Qualitätsaspekte der Tests aus?
    \item In welchem Verhältnis steht der Aufwand für die Erstellung des Kontexts zum erreichten Qualitätsgewinn?
\end{enumerate}

Ergänzend dazu wird untersucht, ob sich die Grenzen eines einmalig bereitgestellten Kontexts überwinden lassen, wenn das Modell Rückmeldung aus der Ausführung seiner eigenen Tests erhält und diese schrittweise korrigiert. Diese Fragestellung geht über den Kernvergleich hinaus und wird daher als Bonus-Untersuchung behandelt. Zur Einordnung der Übertragbarkeit wird der Generierungsablauf zusätzlich mit einem zweiten \ac{LLM} wiederholt. Prompting-Strategien und ein Vergleich der Modellfähigkeiten sind damit nicht Teil dieser Arbeit.

\subsection{Aufbau der Arbeit}

Kapitel~2 erklärt zunächst die Grundlagen zu \acp{LLM}, \ac{E2E}-Testautomatisierung mit Playwright und \ac{OPT}, geht auf die besondere Testkomplexität von WebGIS-Anwendungen ein und ordnet die Arbeit in den aktuellen Stand der Forschung ein. Die Methodik mit den untersuchten Kontextstufen und dem Evaluationsverfahren beschreibt Kapitel~3. Kapitel~4 stellt die dafür implementierte Evaluationsumgebung vor. Kapitel~5 dokumentiert die Durchführung und präsentiert die Ergebnisse des Experiments, die Kapitel~6 anschließend diskutiert. Kapitel~7 schließt die Arbeit mit einem Fazit und Ausblick ab.